\documentclass[11pt]{article}
\usepackage{amsmath,amssymb,amsthm,geometry,hyperref}
\geometry{margin=1in}
\newtheorem{proposition}{Proposition}
\newcommand{\C}{\mathbb C}
\newcommand{\cP}{\mathcal P}
\newcommand{\diag}{\operatorname{diag}}
\newcommand{\tr}{\operatorname{tr}}
\title{The exact translation defect of the transported Fable map}
\author{Independent programme calculation}
\date{20 September 2026}
\begin{document}
\maketitle

The four source points and the polynomial $P_i$ are the exact original
ones in the ES--Fable reader, theorem
\texttt{explicit-four-time-Jacobian-collision}. Their entire polynomial
formulas are reproduced in
\texttt{FOUR\_LABELS\_AND\_CONDUCTOR\_ORDER.tex}. The original
operator and norms are WCF1--6. The four-plane receiving maps and
their full Gram matrices are R1--8 of
\texttt{ACTUAL\_FOUR\_PLANE\_METRIC\_AND\_CONJUGACY.tex}.
This calculation tests an additional structure on exactly those maps:
ordinary polynomial translation. It neither changes the conductor
nor withdraws the established nonlinear conjugacy.

\section{The derivative on the original quotient and target}
Retain $U=\cP_{v+3}/\cP_{v-1}$, $W=\cP_3$, the actual moment
order $v\ge0$, and $T=T_A|_U:U\to W$. Differentiation descends
to the quotient because it preserves $\cP_{v-1}$. Its source
and target matrices in the original ascending monomial frames are
\begin{equation}
 N_v=\begin{pmatrix}0&v+1&0&0\\0&0&v+2&0\\0&0&0&v+3\\0&0&0&0\end{pmatrix},
 \qquad
 N_0=\begin{pmatrix}0&1&0&0\\0&0&2&0\\0&0&0&3\\0&0&0&0\end{pmatrix}.
 \tag{TD1}
\end{equation}
The first column of $N_v$ is zero because the derivative of $S^v$
belongs to the removed kernel. All three superdiagonal entries are
nonzero, so $N_v^4=0$ and
$N_v^3e_3=(v+3)(v+2)(v+1)e_0\ne0$. Both matrices are regular
nilpotent, of rank three.

Every original shift $P(S)\mapsto P(S'+\beta_{8,u,w})$ commutes
with differentiation, so the complete original operator does also:
\begin{equation}
 T\partial_U=\partial_W T.\tag{TD2}
\end{equation}
On $U$ and $W$, exponentiation is a finite cubic sum. It is exactly
the translation $P(S)\mapsto P(S+t)$ on polynomials and its induced
quotient action. No assertion of unitarity in the original Gamma
norm is needed or made.

Retain the ordered root Vandermonde $V_{ja}=\rho_a^j$, the
coefficient matrix $C=(V^{\mathsf T})^{-1}$ of the original
Lagrange polynomials, and the original four-point matrix $K$.
The original conductor coefficient matrix on $U$ is
\[
 (B_v)_{ij}=\begin{cases}\binom{v+j}{i}\mu_{v+j-i},&i\le j,\\
 0,&i>j.\end{cases}
\]
Its invertibility was proved in R1--2. Direct computation gives
$N_0B_v=B_vN_v$: for a nonzero candidate entry this is the binomial
identity
\[
 (i+1)\binom{v+j}{i+1}=(v+j)\binom{v+j-1}{i},
\]
with the same moment $\mu_{v+j-i-1}$ on both sides. The column
$j=0$ is zero on both sides by the actual order $v$.

Let $\Psi=\mathcal YK^{-1}:\C^4\to W$, where
$\mathcal Ye_a=y_a$; its target coefficient matrix is $B_vCK^{-1}$.
The exact pulled-back translation generator is therefore
\begin{equation}
 M=\Psi^{-1}\partial_W\Psi
   =K C^{-1}B_v^{-1}N_0B_v C K^{-1}
   =K V^{\mathsf T}N_v(V^{\mathsf T})^{-1}K^{-1}.
 \tag{TD3}
\end{equation}
This cancellation removes no original moment from $T$: it uses the
proved intertwining identity TD2 to compute the action induced by
that same invertible $T$. In particular $M$ is regular nilpotent,
$M^4=0$ and $M^3\ne0$.

There is a pointwise formula retaining the quotient correction.
On cubic polynomials define
\[
 D_vf(S)=f'(S)+v\frac{f(S)-f(0)}{S}.
\]
The apparent division is exact polynomial division because the
numerator vanishes at zero. Differentiating $S^vf$ and discarding
only the term $vf(0)S^{v-1}$ proves that $D_v$ represents
$\partial_U$ in this cubic model. Consequently
\begin{equation}
 M=KR_vK^{-1},\qquad
 (R_v)_{ba}=\ell_a'(\rho_b)
        +v\frac{\delta_{ba}-\ell_a(0)}{\rho_b}.\tag{TD4}
\end{equation}
The original roots are nonzero, so all divisions displayed in TD4
are valid. In terms of the full original polynomial $h$,
\[
 \ell_a(0)=-\frac{h(0)}{\rho_a h'(\rho_a)},\qquad
 \ell_a'(\rho_b)=\begin{cases}
 \dfrac{h'(\rho_b)}{(\rho_b-\rho_a)h'(\rho_a)},&b\ne a,\\[6pt]
 \displaystyle\sum_{c\ne a}\dfrac1{\rho_a-\rho_c},&b=a.
 \end{cases}
\]
The off-diagonal identity follows by differentiating
$h(S)/((S-\rho_a)h'(\rho_a))$ at its distinct root $\rho_b$;
the diagonal identity follows from the product formula for
$\ell_a$. These equations specify every entry of $M$ in the
original coordinates.

\section{The full polynomial defect and its nonzero linear term}
At the origin the original polynomial has
\begin{equation}
 P_i(0)=0,\qquad J_0=DP_i(0)=\diag(-i,1,2i,-1).\tag{TD5}
\end{equation}
Each diagonal entry is the coefficient of the corresponding original
linear monomial; all other terms of $P_i$ have degree at least two.
For $q=(a,y,z,w)^{\mathsf T}$ put $\xi=Mq$. The entire
infinitesimal translation defect is the polynomial vector
\begin{equation}
 \mathcal D(q)=J(q)\xi-MP_i(q),\qquad J(q)=DP_i(q).\tag{TD6}
\end{equation}
For explicitness, the four rows of the complete Jacobian $J(q)$ are
\begin{align*}
J_{0,*}={}&\bigl(3a^2z+4ay-i,\ 2a^2,\ a^3,\ 0\bigr),\\
J_{2,*}={}&\bigl(
 -3a^2y^2z-4iayz+2aw-4ay^3-10iy^2+3z,\\
 &\hspace{18mm}-2a^3yz-2ia^2z-6a^2y^2-20iay+1,\
 -a^3y^2-2ia^2y+3a,\ a^2\bigr),\\
J_{3,*}={}&\bigl(
 6a^2y^3z+12iay^2z+4awy+8ay^4+2iw-4iy^3-2yz,\\
 &\hspace{18mm}6a^3y^2z+12ia^2yz+2a^2w+16a^2y^3
                  -12iay^2-2az+14y,\\
 &\hspace{18mm}2a^3y^3+6ia^2y^2-2ay+2i,\ 2a^2y+2ia\bigr),\\
J_{4,*}={}&\bigl(
 6a^2y^4z+16iay^3z+2awy^2+8ay^5+2iwy+7iy^4-10y^2z,\\
 &\hspace{18mm}8a^3y^3z+24ia^2y^2z+2a^2wy+20a^2y^4
        +2iaw+28iay^3-20ayz-4iz-9y^2,\\
 &\hspace{18mm}2a^3y^4+8ia^2y^3-10ay^2-4iy,\
 a^2y^2+2iay-1\bigr).
\end{align*}
The row labels $0,2,3,4$ are the original four output labels, not
renumbered coefficient directions. Each entry is obtained by
differentiating the corresponding original monomial. Equations
TD3--6 therefore compute the complete defect, including every
nonlinear term, rather than only its linearization.

Its linear part is
\begin{equation}
 \mathcal D(q)=(J_0M-MJ_0)q+\text{terms of degree at least two}.
 \tag{TD7}
\end{equation}
The displayed commutator is nonzero. Indeed, if $J_0M=MJ_0$,
then for every pair of indices the equality of entries gives
$(\lambda_r-\lambda_s)M_{rs}=0$, where
$(\lambda_1,\lambda_2,\lambda_3,\lambda_4)=(-i,1,2i,-1)$.
These eigenvalues are pairwise distinct, so all off-diagonal entries
of $M$ would vanish. Its nilpotence would then force all diagonal
entries to vanish as well, making $M=0$, contrary to TD3. Thus
$\mathcal D$ is a nonzero, explicitly specified polynomial.

\section{Every nonzero translation and the unchanged original norm}
For any complex translation parameter $t$, define the finite defect
\[
 F_t(q)=P_i(e^{tM}q)-e^{tM}P_i(q),\qquad
 e^{tM}=I+tM+\tfrac12t^2M^2+\tfrac16t^3M^3.
\]
Differentiation in $t$ gives
$\partial_tF_t=MF_t+\mathcal D(e^{tM}q)$ and $F_0=0$.
Multiplying by $e^{-tM}$ and integrating proves the exact formula
\begin{equation}
 F_t(q)=\int_0^t e^{(t-s)M}\mathcal D(e^{sM}q)\,ds.\tag{TD8}
\end{equation}
All expressions in this integral are polynomials; it denotes their
unique polynomial antiderivative with zero value at $t=0$. No
analytic convergence assumption is involved.

For every $t\ne0$, $F_t$ is not the zero polynomial. If it were,
differentiating with respect to $q$ at zero would give
$J_0e^{tM}=e^{tM}J_0$. Put $R=e^{tM}-I$. Then $R^4=0$ and
\[
 tM=\log(I+R)=R-\tfrac12R^2+\tfrac13R^3.
\]
The last equality follows by substituting the cubic exponential and
using $M^4=0$. Commutation with $e^{tM}$ would imply commutation
with $R$, then with this logarithm, and finally with $M$, contrary
to TD7. Thus the noncommutation holds at every nonzero translation
parameter, not only as an infinitesimal statement.

Return now to $W$ with its original Gamma norm and the transported
polynomial $P_W=\Psi P_i\Psi^{-1}$. TD3 and direct substitution give
the complete target defect
\begin{equation}
 P_W(e^{t\partial_W}w)-e^{t\partial_W}P_W(w)
 =\Psi F_t(\Psi^{-1}w).\tag{TD9}
\end{equation}
This is the exact morphism relating the computed defect to the
original polynomial-translation structure in the programme.

Let $G_y$ be the original Gram matrix of the four $y_a$ and put
\[
 G=(K^{-1})^*G_yK^{-1}>0,\qquad C_0=J_0M-MJ_0.
\]
This $G$ is the exact pullback of the original target norm under
$\Psi$. The squared Hilbert--Schmidt norm of the linear part of
the defect, with that same norm on input and output, is
\begin{equation}
 \|C_0\|_{\mathrm{HS},G}^2
 =\tr(G^{-1}C_0^*GC_0)>0.\tag{TD10}
\end{equation}
For the equality, conjugate $C_0$ by the positive square root
$G^{1/2}$ and take the sum of the squared moduli of all entries.
Cyclicity of the finite matrix trace gives TD10. Positivity is
strict because $C_0\ne0$. The source pullback remains
$(K^{-1})^*G_xK^{-1}$, as in R3--5; the original conductor maps
that source to this target and still obeys its exact nonlinear
conjugacy square. None of the preceding expressions equates these
two metrics or discards their masses.

The transported Fable map therefore has a completely proved relation
to the original conductor and a completely proved, nonzero defect
with respect to its translation symmetry. Conjugacy supplies the
former relation; it does not silently supply translation equivariance
or a bound on the original remaining programme terms.
\end{document}
